\documentclass[11pt,a4paper]{article}

% --- Packages ---
\usepackage[margin=1in]{geometry}
\usepackage{amsmath,amssymb,amsthm,mathtools}
\usepackage{thmtools}
\usepackage{enumitem}
\usepackage{hyperref}
\usepackage[dvipsnames]{xcolor}
\usepackage{tikz}
\usetikzlibrary{arrows.meta,positioning,calc,decorations.pathmorphing,shapes.geometric}
\usepackage[ruled,vlined,linesnumbered]{algorithm2e}
\usepackage{bm}
\usepackage{nicefrac}
\usepackage{microtype}
\usepackage{booktabs}
\usepackage{cleveref}

\hypersetup{
  colorlinks=true,
  linkcolor=NavyBlue,
  citecolor=OliveGreen,
  urlcolor=BrickRed
}

% --- Theorem environments ---
\theoremstyle{plain}
\newtheorem{theorem}{Theorem}[section]
\newtheorem{proposition}[theorem]{Proposition}
\newtheorem{lemma}[theorem]{Lemma}
\newtheorem{corollary}[theorem]{Corollary}

\theoremstyle{definition}
\newtheorem{definition}[theorem]{Definition}
\newtheorem{assumption}[theorem]{Assumption}
\newtheorem{example}[theorem]{Example}

\theoremstyle{remark}
\newtheorem{remark}[theorem]{Remark}

% --- Notation shortcuts ---
\newcommand{\R}{\mathbb{R}}
\newcommand{\E}{\mathbb{E}}
\newcommand{\Prob}{\mathbb{P}}
\newcommand{\cF}{\mathcal{F}}
\newcommand{\cH}{\mathcal{H}}
\newcommand{\cS}{\mathcal{S}}
\newcommand{\cA}{\mathcal{A}}
\newcommand{\cV}{\mathcal{V}}
\newcommand{\cX}{\mathcal{X}}
\newcommand{\cZ}{\mathcal{Z}}
\newcommand{\cT}{\mathcal{T}}
\newcommand{\cM}{\mathcal{M}}
\newcommand{\cL}{\mathcal{L}}
\newcommand{\cD}{\mathcal{D}}
\newcommand{\cP}{\mathcal{P}}
\newcommand{\cN}{\mathcal{N}}
\newcommand{\cC}{\mathcal{C}}
\newcommand{\cB}{\mathcal{B}}
\newcommand{\cW}{\mathcal{W}}
\newcommand{\KL}{\mathrm{KL}}
\newcommand{\TV}{\mathrm{TV}}
\newcommand{\tr}{\mathrm{tr}}
\newcommand{\diag}{\mathrm{diag}}
\newcommand{\softmax}{\mathrm{softmax}}
\newcommand{\Attn}{\mathrm{Attn}}
\newcommand{\vocab}{\mathcal{V}}
\newcommand{\ind}{\mathbf{1}}
\DeclareMathOperator*{\argmax}{arg\,max}
\DeclareMathOperator*{\argmin}{arg\,min}

\title{%
  \textbf{Autoregressive Language Models as\\
  Bayesian Optimal Controllers over\\
  Hidden Markovian State Spaces}\\[0.6em]
  \large A Unified Theoretical Framework
}

\author{%
  A Theoretical Synthesis
}

\date{April 2026}

\begin{document}
\maketitle

\begin{abstract}
We develop a unified mathematical framework for understanding large autoregressive
transformer language models through the lens of three interconnected formalisms:
\emph{hidden Markov processes}, \emph{stochastic optimal control}, and
\emph{Bayesian inference}. We formalise the claim that the internal hidden state
of a transformer defines a high-dimensional, continuous-state hidden Markov
process whose effective memory decays exponentially, yielding asymptotic
memorylessness governed by the spectral properties of an induced transfer
operator. We then show that autoregressive token generation can be cast as a
stochastic optimal control problem in which the model's hidden state serves as
the control state, token emissions are control actions, and the learned value
function encodes a Bellman-optimal policy over sequence trajectories. Finally, we
demonstrate that the pretrained weights define a prior measure over token
sequences, reinforcement learning from human feedback (RLHF) reshapes this prior
into a value-weighted posterior, and in-context conditioning performs amortised
variational Bayesian updating within each conversation. The three perspectives
are proven to be mutually consistent, yielding a coherent theoretical picture in
which memory decay, trajectory optimisation, and posterior concentration are
manifestations of a single underlying generative mechanism.
\end{abstract}

\tableofcontents
\newpage

%=============================================================================
\section{Introduction and Motivation}\label{sec:intro}
%=============================================================================

Modern large language models (LLMs) based on the transformer architecture
generate text by iteratively sampling tokens from a conditional distribution.
While the engineering details of these models are well understood, a unified
\emph{theoretical} framework that explains their behaviour through the lens of
classical probability, control theory, and Bayesian statistics remains
fragmented.

In this work, we propose that three classical formalisms---hidden Markov
processes, stochastic optimal control, and Bayesian inference---each capture a
complementary facet of transformer behaviour, and that together they form a
coherent and mutually reinforcing picture.

\paragraph{Core thesis.}
Let $\bm{h}_t \in \R^d$ denote the hidden representation of a transformer at
decoding step $t$, and let $x_t \in \vocab$ denote the emitted token. We argue:
\begin{enumerate}[label=(\roman*)]
  \item The stochastic process $\{(\bm{h}_t, x_t)\}_{t \geq 1}$ constitutes a
        \emph{hidden Markov process} on a continuous state space with
        exponentially decaying memory.
  \item Autoregressive generation solves a \emph{stochastic optimal control
        problem} in which $\bm{h}_t$ is the state, $x_t$ is the action, and the
        model has internalised an approximate value function.
  \item The model parameters encode a \emph{Bayesian prior} over sequences that
        is updated in-context via amortised approximate inference.
\end{enumerate}

\paragraph{Notation.} Throughout, we use $\vocab$ for the token vocabulary
(finite, $|\vocab| = V$), $d$ for the hidden dimension, $T$ for the sequence
length or horizon, $\bm{x}_{1:t} = (x_1, \ldots, x_t)$ for token prefixes, and
$\bm{h}_t \in \R^d$ for the hidden state at step $t$. Probability measures are
denoted $\Prob$, densities $p$, and expectations $\E$. Parameters of the model
are collectively denoted $\theta \in \Theta$.

%=============================================================================
\section{Preliminaries}\label{sec:prelim}
%=============================================================================

\subsection{Markov Chains on General State Spaces}

\begin{definition}[Markov kernel]
Let $(\cS, \cB(\cS))$ be a measurable space. A \emph{Markov kernel} (or
transition kernel) is a map $\kappa: \cS \times \cB(\cS) \to [0,1]$ such that
(i) for each $s \in \cS$, $\kappa(s, \cdot)$ is a probability measure on
$\cB(\cS)$; and (ii) for each $B \in \cB(\cS)$, $\kappa(\cdot, B)$ is
measurable.
\end{definition}

\begin{definition}[Markov chain on a general state space]
A sequence of random variables $\{S_t\}_{t \geq 0}$ taking values in $\cS$ is a
\emph{Markov chain} with kernel $\kappa$ if for all $t$ and all
$B \in \cB(\cS)$,
\[
  \Prob(S_{t+1} \in B \mid S_0, \ldots, S_t) = \kappa(S_t, B) \quad \text{a.s.}
\]
\end{definition}

\begin{definition}[Ergodicity and mixing]\label{def:ergodic}
A Markov chain with kernel $\kappa$ and stationary distribution $\pi$ is
\emph{geometrically ergodic} if there exists $\rho \in (0,1)$ and a measurable
function $M: \cS \to [0,\infty)$ such that for all $s \in \cS$ and $t \geq 0$,
\[
  \|\kappa^t(s, \cdot) - \pi(\cdot)\|_{\TV} \leq M(s)\, \rho^t,
\]
where $\|\cdot\|_{\TV}$ denotes the total variation norm and $\kappa^t$ is the
$t$-fold composition of $\kappa$.
\end{definition}

The quantity $\tau_{\mathrm{mix}} = \lceil -1/\log \rho \rceil$ is the
\emph{mixing time}: beyond this horizon, the chain's distribution is
$\varepsilon$-close to $\pi$ regardless of the initial state.

\subsection{Hidden Markov Models}

\begin{definition}[Hidden Markov Model]
A \emph{Hidden Markov Model} (HMM) is a tuple $(\cS, \cX, \kappa, g, \nu)$
where:
\begin{itemize}[nosep]
  \item $\cS$ is the (possibly continuous) hidden state space,
  \item $\cX$ is the observation space,
  \item $\kappa$ is a Markov transition kernel on $\cS$,
  \item $g: \cS \to \cP(\cX)$ is the emission distribution,
  \item $\nu \in \cP(\cS)$ is the initial distribution.
\end{itemize}
The generative model factorises as
\[
  s_0 \sim \nu, \qquad
  s_{t+1} \sim \kappa(s_t, \cdot), \qquad
  x_t \sim g(s_t).
\]
\end{definition}

\subsection{Stochastic Optimal Control}

\begin{definition}[Discrete-time stochastic control problem]
A \emph{stochastic optimal control problem} is specified by:
\begin{itemize}[nosep]
  \item A state space $\cS$, action space $\cA$,
  \item State dynamics $s_{t+1} \sim f(s_t, a_t, \xi_t)$ where
        $\xi_t$ is i.i.d.\ noise,
  \item A per-step reward $r: \cS \times \cA \to \R$,
  \item A terminal reward $R_T: \cS \to \R$,
  \item A discount factor $\gamma \in (0,1]$.
\end{itemize}
The objective is to find a policy $\pi: \cS \to \cP(\cA)$ maximising
\[
  J(\pi) = \E^{\pi}\!\left[\sum_{t=0}^{T} \gamma^t\, r(s_t, a_t) + \gamma^T R_T(s_T)\right].
\]
\end{definition}

\begin{theorem}[Bellman optimality, discrete-time]\label{thm:bellman}
The optimal value function $V^*: \cS \to \R$ satisfies
\[
  V^*(s) = \max_{a \in \cA}\!\left\{r(s,a) + \gamma\, \E_{s' \sim f(s,a,\cdot)}[V^*(s')]\right\}
\]
and the optimal policy is $\pi^*(s) = \argmax_{a}\{r(s,a) + \gamma\,\E[V^*(f(s,a,\xi))]\}$.
\end{theorem}

\subsection{Bayesian Inference}

\begin{definition}[Bayesian framework]
Given a prior $p(\theta)$ over parameters and a likelihood $p(\cD|\theta)$ for
data $\cD$, the posterior is
\[
  p(\theta|\cD) = \frac{p(\cD|\theta)\,p(\theta)}{p(\cD)}, \qquad
  p(\cD) = \int p(\cD|\theta)\,p(\theta)\,d\theta.
\]
\end{definition}

In the sequel, we will use the Bayesian framework not over model parameters
$\theta$ (which are fixed at inference time) but over \emph{sequence
continuations}: the model parameters define a prior over token sequences, and
the observed context serves as evidence.


%=============================================================================
\section{The Transformer as a Hidden Markov Process}\label{sec:hmm}
%=============================================================================

\subsection{Defining the Hidden State}

Consider a transformer with parameters $\theta$ operating autoregressively.
At decoding step $t$, the model has consumed tokens $\bm{x}_{1:t}$ and
produced an internal representation that we denote $\bm{h}_t \in \R^d$.
Formally, $\bm{h}_t$ comprises the residual stream activations and, in
cached inference, the key-value (KV) cache:
\[
  \bm{h}_t = \Phi_\theta(\bm{x}_{1:t}),
\]
where $\Phi_\theta: \vocab^{\leq T} \to \R^d$ is the deterministic encoding
map defined by the transformer's forward pass.

\begin{proposition}[Markov property of the hidden state]\label{prop:markov}
The next-token distribution depends on $\bm{x}_{1:t}$ \emph{only through}
$\bm{h}_t$:
\[
  p_\theta(x_{t+1} \mid \bm{x}_{1:t})
  = p_\theta(x_{t+1} \mid \bm{h}_t)
  = \softmax\!\bigl(W_{\mathrm{out}}\, \bm{h}_t^{(L)}\bigr)_{x_{t+1}},
\]
where $\bm{h}_t^{(L)}$ is the final-layer representation and $W_{\mathrm{out}}
\in \R^{V \times d}$ is the unembedding matrix. Hence the pair process
$\{(\bm{h}_t, x_t)\}_{t \geq 1}$ is Markovian with transition kernel
\begin{align}
  \kappa_\theta\bigl((\bm{h}_t, x_t),\, B \times \{v\}\bigr)
  &= p_\theta(v \mid \bm{h}_t)\cdot
     \ind\!\bigl[\Phi_\theta(\bm{x}_{1:t}, v) \in B\bigr]
     \label{eq:kernel}
\end{align}
for $B \in \cB(\R^d)$ and $v \in \vocab$.
\end{proposition}

\begin{proof}
Since $\bm{h}_{t+1} = \Phi_\theta(\bm{x}_{1:t}, x_{t+1})$ is a deterministic
function of $(\bm{h}_t, x_{t+1})$ (via the incremental KV-cache update), and
$x_{t+1}$ is sampled from $p_\theta(\cdot \mid \bm{h}_t)$, the state at
$t+1$ depends on the past only through $(\bm{h}_t, x_t)$. The emission
$x_{t+1}$ depends only on $\bm{h}_t$. This establishes the Markov property.
\end{proof}

\begin{remark}[Non-trivial Markov structure]
One might object that this is ``trivially'' Markovian since $\bm{h}_t$
deterministically encodes the full history $\bm{x}_{1:t}$. However, the
claim becomes substantive when we consider the \emph{effective} memory of
the process---i.e., the degree to which distant tokens actually influence
future predictions through the hidden state.
\end{remark}

\subsection{The HMM Interpretation}

\begin{definition}[Transformer HMM]\label{def:thmm}
The \emph{Transformer Hidden Markov Model} (T-HMM) is the tuple
$(\cS, \vocab, \kappa_\theta, g_\theta, \nu_\theta)$ where:
\begin{itemize}[nosep]
  \item $\cS = \R^d$ is the continuous hidden state space,
  \item $\vocab$ is the finite token vocabulary (observation space),
  \item $\kappa_\theta$ is the kernel from~\eqref{eq:kernel},
  \item $g_\theta(\bm{h}) = \softmax(W_{\mathrm{out}}\, \bm{h}^{(L)})$ is the
        emission distribution,
  \item $\nu_\theta$ is the distribution over initial hidden states induced by
        the start-of-sequence token.
\end{itemize}
\end{definition}

Unlike classical HMMs with finite discrete hidden states, the T-HMM has a
continuous, high-dimensional hidden space with \emph{deterministic} state
transitions conditioned on the emitted token. This places it in the family of
\emph{deterministic-transition HMMs}, closely related to predictive state
representations and observable operator models.

\subsection{Exponential Decay of Influence}\label{sec:decay}

The central empirical and theoretical observation is that the influence of
token $x_s$ on the prediction at step $t \gg s$ decays with the lag $t - s$.

\begin{definition}[Influence function]
The \emph{influence} of token position $s$ on the prediction at position $t$ is
\[
  I(s, t) \triangleq
  \E\!\left[\KL\!\Big(
    p_\theta(\cdot \mid \bm{x}_{1:t})
    \;\Big\|\;
    p_\theta(\cdot \mid \bm{x}_{1:t}^{\setminus s})
  \Big)\right],
\]
where $\bm{x}_{1:t}^{\setminus s}$ denotes the sequence with token $x_s$
replaced by a draw from the marginal $p_\theta(x_s)$.
\end{definition}

\begin{theorem}[Exponential influence decay]\label{thm:decay}
Suppose the hidden-state transition map $F_\theta: \R^d \times \vocab \to \R^d$
defined by $\bm{h}_{t+1} = F_\theta(\bm{h}_t, x_{t+1})$ is Lipschitz in its
first argument with constant $\lambda < 1$ uniformly over $\vocab$:
\[
  \|F_\theta(\bm{h}, v) - F_\theta(\bm{h}', v)\| \leq \lambda\, \|\bm{h} - \bm{h}'\|
  \quad \forall\, v \in \vocab,\; \bm{h}, \bm{h}' \in \R^d.
\]
Then the influence decays geometrically:
\[
  I(s, t) \leq C\, \lambda^{2(t-s)}
\]
for a constant $C$ depending on the Lipschitz constant of $g_\theta$ and the
diameter of the hidden state space.
\end{theorem}

\begin{proof}[Proof sketch]
A perturbation $\delta \bm{h}_s = \bm{h}_s - \bm{h}_s'$ at step $s$ propagates
forward as
\[
  \|\delta \bm{h}_t\| \leq \lambda^{t-s}\, \|\delta \bm{h}_s\|
\]
by iterated application of the contraction. Since $g_\theta$ is Lipschitz (the
softmax composed with a linear map), the KL divergence between the resulting
emission distributions is bounded by $O(\|\delta \bm{h}_t\|^2)$ via Pinsker's
inequality and a local quadratic bound on KL, giving the $\lambda^{2(t-s)}$
rate. \qed
\end{proof}

\begin{remark}[Connection to mixing]
In the language of \Cref{def:ergodic}, the T-HMM is geometrically ergodic
with rate $\rho = \lambda^2$ whenever the contraction condition holds. The
mixing time $\tau_{\mathrm{mix}} = O(1/|\log \lambda|)$ determines the
effective memory horizon of the process.
\end{remark}

\begin{remark}[Attention as non-uniform contraction]
In practice, multi-head attention can selectively break the contraction for
specific positions (e.g., syntactic anchors, proper nouns), creating
\emph{non-uniform} decay. The process is better modelled as a
\emph{time-inhomogeneous} Markov chain with position-dependent contraction
rates $\lambda_t$. The cumulative decay over a window $[s, t]$ is then
$\prod_{\tau=s}^{t-1} \lambda_\tau$, which is still exponentially small when
$\bar{\lambda} = (\prod \lambda_\tau)^{1/(t-s)} < 1$ on average.
\end{remark}

\subsection{Long-Run Memorylessness}

\begin{corollary}[Asymptotic memorylessness]\label{cor:memoryless}
Under the contraction condition of \Cref{thm:decay}, for any $\varepsilon > 0$
there exists $\tau(\varepsilon)$ such that for all $t - s > \tau(\varepsilon)$,
\[
  I(s, t) < \varepsilon.
\]
In the limit $t \to \infty$ (with stationarity), the process is \emph{memoryless
in the long run}: the predictive distribution converges to a fixed point
independent of arbitrarily distant history.
\end{corollary}

This formalises the intuition: at every step, the process has \emph{local
memory}---the current state $\bm{h}_t$ fully determines the next transition.
But the influence of any particular historical token on that state is
\emph{vanishingly small} beyond the mixing horizon. The process is neither
memoryless in the white-noise sense (where the present carries no information)
nor long-memory in the fractional Brownian motion sense. It occupies the
intermediate regime of \emph{exponentially decaying, finite-horizon memory}.


%=============================================================================
\section{Autoregressive Generation as Optimal Control}\label{sec:control}
%=============================================================================

We now reinterpret the generative process through the lens of stochastic
optimal control.

\subsection{The Control Problem Formulation}

\begin{definition}[Autoregressive control problem]\label{def:arp}
Define the \emph{autoregressive control problem} (ACP) as follows:
\begin{itemize}[nosep]
  \item \textbf{State space:} $\cS = \R^d$ (hidden representations).
  \item \textbf{Action space:} $\cA = \vocab$ (token vocabulary).
  \item \textbf{Dynamics:} $\bm{h}_{t+1} = F_\theta(\bm{h}_t, x_{t+1})$
        (deterministic given action).
  \item \textbf{Policy:} $\pi_\theta(\cdot \mid \bm{h}_t) = g_\theta(\bm{h}_t)
        \in \cP(\vocab)$.
  \item \textbf{Reward:} To be specified by the training objective.
  \item \textbf{Horizon:} $T$ (sequence length or end-of-sequence).
\end{itemize}
\end{definition}

\subsection{Pretraining as Maximum Likelihood Control}

Under maximum likelihood training, the implicit reward is the log-probability of
the observed next token:
\[
  r_{\mathrm{MLE}}(\bm{h}_t, x_{t+1}) = \log p_\theta(x_{t+1} \mid \bm{h}_t).
\]
The pretraining objective maximises the cumulative reward over the data
distribution:
\[
  J_{\mathrm{MLE}}(\theta)
  = \E_{\bm{x} \sim \cD}\!\left[\sum_{t=0}^{T-1} \log p_\theta(x_{t+1} \mid \bm{h}_t)\right]
  = \E_{\bm{x} \sim \cD}[\log p_\theta(\bm{x})].
\]

\begin{proposition}[MLE as undiscounted control]
Maximum likelihood pretraining is equivalent to solving the ACP with
$\gamma = 1$, $r(\bm{h}_t, a) = \log p_\theta(a \mid \bm{h}_t)$, and $R_T = 0$,
under the constraint that the policy matches the model's own conditional
distribution.
\end{proposition}

This is a \emph{self-consistent} control problem: the reward function
\emph{is} the policy's own log-probability, and training adjusts $\theta$
until the policy is optimal under its own induced reward.

\subsection{RLHF as Reward Shaping}

Reinforcement Learning from Human Feedback introduces an external reward
signal. Let $R_\phi: \vocab^T \to \R$ denote a learned reward model with
parameters $\phi$, trained on human preference data.

\begin{definition}[RLHF objective]
The RLHF objective is
\begin{equation}\label{eq:rlhf_obj}
  J_{\mathrm{RLHF}}(\theta)
  = \E_{\bm{x} \sim \pi_\theta}\!\big[R_\phi(\bm{x})\big]
    - \beta\, \KL\!\big(\pi_\theta \| \pi_{\mathrm{ref}}\big),
\end{equation}
where $\pi_{\mathrm{ref}}$ is the pretrained (reference) policy and
$\beta > 0$ controls the strength of the KL penalty.
\end{definition}

\begin{theorem}[Optimal RLHF policy]\label{thm:rlhf_optimal}
The policy maximising~\eqref{eq:rlhf_obj} has the Gibbs form
\begin{equation}\label{eq:gibbs}
  \pi^*_\theta(\bm{x}) \propto \pi_{\mathrm{ref}}(\bm{x})\,
  \exp\!\bigl(\beta^{-1} R_\phi(\bm{x})\bigr).
\end{equation}
\end{theorem}

\begin{proof}
This is the standard result from the variational characterisation of the
KL-regularised optimisation. The functional
$\cL[\pi] = \E_\pi[R_\phi] - \beta\,\KL(\pi\|\pi_{\mathrm{ref}})$
is strictly concave in $\pi$ and its unique maximiser is the tilted
distribution~\eqref{eq:gibbs}. \qed
\end{proof}

\begin{remark}[Control interpretation]
In the control framework, RLHF replaces the self-referential MLE reward with a
\emph{trajectory-level} reward $R_\phi(\bm{x}_{1:T})$, and the KL penalty
acts as a \emph{control cost} penalising deviation from the reference dynamics.
The optimal policy~\eqref{eq:gibbs} is precisely the solution to the
KL-regularised stochastic control problem, a well-known result in
linearly-solvable MDPs and path-integral control.
\end{remark}

\subsection{The Partition Function and the Free-Boundary Bellman Equation}

The key insight is that RLHF rewards \emph{entire sequences}, not individual
tokens. A model that must emit tokens sequentially but is evaluated on
trajectory-level reward is \emph{forced} to decompose that trajectory reward
into per-step decisions. Moreover, the horizon $T$ is not fixed: the model
decides \emph{when to stop} by emitting a distinguished end-of-sequence token
$\langle\mathrm{eos}\rangle \in \vocab$. This makes $T$ endogenous---the
stopping time is part of the solution, not a given.

\begin{definition}[Stopping time]
The \emph{generation stopping time} is the random variable
\[
  \tau = \inf\{t \geq 1 : x_t = \langle\mathrm{eos}\rangle\},
\]
i.e., the first time the model emits the end-of-sequence token. The
completed sequence is $\bm{x}_{1:\tau}$, and the trajectory reward is
$R_\phi(\bm{x}_{1:\tau})$.
\end{definition}

\begin{definition}[Trajectory partition function with endogenous stopping]\label{def:partition}
For a hidden state $\bm{h}_t$ at step $t$, define the \emph{trajectory
partition function}
\begin{equation}\label{eq:partition}
  Z(\bm{h}_t) \triangleq
  \E_{\pi_{\mathrm{ref}}}\!\left[
    \exp\!\bigl(\beta^{-1} R_\phi(\bm{x}_{1:\tau})\bigr)
    \;\Big|\; \bm{h}_t
  \right],
\end{equation}
where the expectation is over all continuations under $\pi_{\mathrm{ref}}$,
including the stochastic stopping time $\tau \geq t$.
\end{definition}

\begin{theorem}[Free-boundary backward recursion]\label{thm:Z_recursion}
The partition function satisfies the backward recursion
\begin{equation}\label{eq:Z_backward}
  Z(\bm{h}_t)
  = \underbrace{\pi_{\mathrm{ref}}(\langle\mathrm{eos}\rangle \mid \bm{h}_t)
    \cdot \exp\!\bigl(\beta^{-1} R_\phi(\bm{x}_{1:t})\bigr)}_{\displaystyle
    S(\bm{h}_t)\;\text{: \emph{stopping value}}}
  + \underbrace{\sum_{v \neq \langle\mathrm{eos}\rangle}
    \pi_{\mathrm{ref}}(v \mid \bm{h}_t)
    \cdot Z\!\bigl(F_\theta(\bm{h}_t, v)\bigr)}_{\displaystyle
    C(\bm{h}_t)\;\text{: \emph{continuation value}}}.
\end{equation}
The stopping value $S(\bm{h}_t)$ is the expected exponentiated reward from
terminating now; the continuation value $C(\bm{h}_t)$ is the expected
exponentiated reward from emitting one more non-terminal token and proceeding
optimally thereafter.
\end{theorem}

\begin{proof}
By the law of total expectation, conditioning on whether the next token is
$\langle\mathrm{eos}\rangle$ or not:
\begin{align}
  Z(\bm{h}_t)
  &= \E_{\pi_{\mathrm{ref}}}\!\left[
       \exp\!\bigl(\beta^{-1} R_\phi(\bm{x}_{1:\tau})\bigr)
       \;\Big|\; \bm{h}_t
     \right] \notag \\
  &= \pi_{\mathrm{ref}}(\langle\mathrm{eos}\rangle \mid \bm{h}_t)
     \cdot \exp\!\bigl(\beta^{-1} R_\phi(\bm{x}_{1:t})\bigr) \notag \\
  &\quad + \sum_{v \neq \langle\mathrm{eos}\rangle}
     \pi_{\mathrm{ref}}(v \mid \bm{h}_t)
     \cdot \E_{\pi_{\mathrm{ref}}}\!\left[
       \exp\!\bigl(\beta^{-1} R_\phi(\bm{x}_{1:\tau})\bigr)
       \;\Big|\; \bm{h}_{t+1} = F_\theta(\bm{h}_t, v)
     \right] \notag \\
  &= S(\bm{h}_t) + C(\bm{h}_t). \tag*{\qed}
\end{align}
\end{proof}

\begin{remark}[The free-boundary structure]
Equation~\eqref{eq:Z_backward} is structurally different from the fixed-horizon
Bellman equation. In the fixed-horizon case, the terminal condition is given
at a known $T$. Here, the ``terminal condition'' $\exp(\beta^{-1}
R_\phi(\bm{x}_{1:t}))$ is available at \emph{every} step---the model can
always choose to stop. The recursion simultaneously determines:
\begin{enumerate}[label=(\alph*),nosep]
  \item The value function $V^*(\bm{h}_t) = \beta \log Z(\bm{h}_t)$.
  \item The optimal stopping region
        $\cS^* = \{\bm{h} : S(\bm{h}) \geq C(\bm{h})\}$---the set of
        hidden states where terminating is preferred to continuing.
\end{enumerate}
The stopping boundary $\partial \cS^*$ is not given a priori; it is a
\emph{free boundary} that emerges from the recursion.
\end{remark}

\begin{definition}[Soft value function]
Define the \emph{soft value function} as
\begin{equation}\label{eq:soft_V}
  V^*(\bm{h}_t) \triangleq \beta \log Z(\bm{h}_t).
\end{equation}
\end{definition}

\begin{theorem}[Free-boundary soft Bellman equation]\label{thm:soft_bellman}
The soft value function satisfies
\begin{equation}\label{eq:soft_bellman}
  \exp\!\bigl(\beta^{-1} V^*(\bm{h}_t)\bigr)
  = \pi_{\mathrm{ref}}(\langle\mathrm{eos}\rangle \mid \bm{h}_t)
    \cdot \exp\!\bigl(\beta^{-1} R_\phi(\bm{x}_{1:t})\bigr)
  + \!\!\sum_{v \neq \langle\mathrm{eos}\rangle}\!\!
    \pi_{\mathrm{ref}}(v \mid \bm{h}_t)
    \cdot \exp\!\bigl(\beta^{-1} V^*\!\bigl(F_\theta(\bm{h}_t, v)\bigr)\bigr).
\end{equation}
This is the \emph{free-boundary} soft Bellman equation: it is the Bellman
equation of a linearly-solvable MDP with an endogenous stopping time.
\end{theorem}

\begin{proof}
Immediate from \Cref{thm:Z_recursion} and \eqref{eq:soft_V}: apply
$\beta \log$ to both sides of~\eqref{eq:Z_backward}. \qed
\end{proof}

\begin{remark}[Connection to optimal stopping theory]
The recursion~\eqref{eq:soft_bellman} is the \emph{Snell envelope} of the
process $\{R_\phi(\bm{x}_{1:t})\}_{t \geq 1}$ under the soft-max aggregation
induced by $\beta$. In the limit $\beta \to 0$ (deterministic policy), it
reduces to the classical optimal stopping rule:
\[
  V^*(\bm{h}_t)
  = \max\!\Big\{
      \underbrace{R_\phi(\bm{x}_{1:t})}_{\text{stop}},\;\;
      \underbrace{\max_{v \neq \langle\mathrm{eos}\rangle}
        V^*\!\bigl(F_\theta(\bm{h}_t, v)\bigr)}_{\text{continue}}
    \Big\},
\]
which is the Bellman equation of an American option: exercise (emit EOS,
collect $R_\phi$) or continue (emit another token, reach a potentially
higher-value state). The optimal stopping boundary $\partial\cS^*$ is
the \emph{exercise boundary} of this option.
\end{remark}

\begin{remark}[Backward induction and lookahead]
The endogenous stopping time forces the model to perform implicit backward
induction. At each step, the model does not ask ``what is the best
complete response to the user's query?''---a planning problem requiring
knowledge of the full trajectory. Instead it asks ``what is the next token
that brings the hidden state closer to the optimal stopping region while
accumulating high value-to-go?'' This is a \emph{local} question answered
by the gradient $\nabla_{\bm{h}} V^*(\bm{h}_t)$: the direction in
hidden-state space along which the value function increases most steeply
toward $\cS^*$. The model surfs the value landscape toward the stopping
boundary, one token at a time.

When this gradient is well-defined and smooth, local hill-climbing in
advantage space produces globally coherent responses. When it is flat or
noisy---the model has poor value estimates in that region of hidden
space---the result is rambling, premature stopping, or inability to
conclude. The model has lost the gradient toward $\partial\cS^*$.
\end{remark}

\begin{theorem}[Optimal autoregressive policy via $Z$]\label{thm:opt_policy_Z}
The optimal policy for the RLHF objective~\eqref{eq:rlhf_obj}, with
endogenous stopping, factorises autoregressively at each step as
\begin{equation}\label{eq:opt_policy_Z}
  \pi^*(v \mid \bm{h}_t)
  = \pi_{\mathrm{ref}}(v \mid \bm{h}_t)
    \cdot \frac{Z\!\bigl(F_\theta(\bm{h}_t, v)\bigr)}{Z(\bm{h}_t)},
    \qquad v \neq \langle\mathrm{eos}\rangle,
\end{equation}
and for the stopping action:
\begin{equation}\label{eq:opt_policy_eos}
  \pi^*(\langle\mathrm{eos}\rangle \mid \bm{h}_t)
  = \pi_{\mathrm{ref}}(\langle\mathrm{eos}\rangle \mid \bm{h}_t)
    \cdot \frac{\exp\!\bigl(\beta^{-1} R_\phi(\bm{x}_{1:t})\bigr)}{Z(\bm{h}_t)}.
\end{equation}
The model emits EOS when $R_\phi(\bm{x}_{1:t})$ (the reward from
stopping) is large relative to $V^*(F_\theta(\bm{h}_t, v))$
(the value from continuing).
\end{theorem}

\begin{proof}
The optimal joint distribution is $\pi^*(\bm{x}_{1:\tau}) \propto
\pi_{\mathrm{ref}}(\bm{x}_{1:\tau}) \exp(\beta^{-1} R_\phi(\bm{x}_{1:\tau}))$
by \Cref{thm:rlhf_optimal}, where $\tau$ is the stopping time. Factoring
autoregressively and using the definition of $Z$:

For $v \neq \langle\mathrm{eos}\rangle$:
\begin{align}
  \pi^*(v \mid \bm{h}_t)
  &= \frac{\sum_{\tau > t} \sum_{\bm{x}_{t+2:\tau}}
     \pi_{\mathrm{ref}}(v \mid \bm{h}_t)
     \pi_{\mathrm{ref}}(\bm{x}_{t+2:\tau} \mid F_\theta(\bm{h}_t, v))
     \exp(\beta^{-1} R_\phi(\bm{x}_{1:\tau}))}{Z(\bm{h}_t)} \notag \\
  &= \pi_{\mathrm{ref}}(v \mid \bm{h}_t) \cdot
     \frac{Z(F_\theta(\bm{h}_t, v))}{Z(\bm{h}_t)}.
\end{align}

For $v = \langle\mathrm{eos}\rangle$ (stopping at step $t+1$, so the
sequence terminates as $\bm{x}_{1:t+1}$ with $x_{t+1} = \langle\mathrm{eos}\rangle$):
\begin{align}
  \pi^*(\langle\mathrm{eos}\rangle \mid \bm{h}_t)
  &= \frac{\pi_{\mathrm{ref}}(\langle\mathrm{eos}\rangle \mid \bm{h}_t)
     \cdot \exp(\beta^{-1} R_\phi(\bm{x}_{1:t}))}{Z(\bm{h}_t)}.
\end{align}
One can verify that these sum to 1 by~\eqref{eq:Z_backward}. \qed
\end{proof}

\begin{remark}[The $Z$-ratio as the central object]
The ratio
\begin{equation}\label{eq:Z_ratio}
  \Gamma(\bm{h}_t, v) \triangleq
  \begin{cases}
    \dfrac{Z\!\bigl(F_\theta(\bm{h}_t, v)\bigr)}{Z(\bm{h}_t)}
    & v \neq \langle\mathrm{eos}\rangle, \\[1em]
    \dfrac{\exp\!\bigl(\beta^{-1} R_\phi(\bm{x}_{1:t})\bigr)}{Z(\bm{h}_t)}
    & v = \langle\mathrm{eos}\rangle,
  \end{cases}
\end{equation}
is the single mathematical object that unifies all three frameworks, as we
shall prove in \Cref{sec:unify}. For continuation tokens, $\Gamma$ measures
the relative value of continuing via $v$. For EOS, $\Gamma$ measures the
relative value of stopping now versus continuing. The model emits EOS
when $\Gamma(\bm{h}_t, \langle\mathrm{eos}\rangle)$ dominates all
$\Gamma(\bm{h}_t, v)$ for $v \neq \langle\mathrm{eos}\rangle$.
\end{remark}

\begin{corollary}[The model must encode the value function]\label{cor:must_encode}
Since the RLHF-trained model approximates $\pi^*$, and $\pi^*$ factorises
through $Z$ via~\eqref{eq:opt_policy_Z}--\eqref{eq:opt_policy_eos}, the
model's hidden state $\bm{h}_t$ must encode sufficient information to
approximate both $V^*(\bm{h}_t) = \beta \log Z(\bm{h}_t)$ (the
continuation value) and $R_\phi(\bm{x}_{1:t})$ (the stopping value).
The residual stream carries the log-partition function of the
value-weighted future \emph{and} the current accumulated reward.
\end{corollary}

\begin{remark}[Chain-of-thought as stopping-time extension]
Chain-of-thought (CoT) prompting delays the stopping decision by emitting
intermediate reasoning tokens. In the free-boundary framework, CoT
shifts the optimal stopping region: the reasoning tokens navigate the
hidden state to a region where $C(\bm{h}_t) \gg S(\bm{h}_t)$ and then
eventually to a region where $S(\bm{h}_t)$ is large and
$R_\phi(\bm{x}_{1:t})$ has been increased by the intervening reasoning.
The model is given more time to climb the value landscape before being
forced to exercise the option.
\end{remark}

\subsection{Pontryagin vs.\ Dynamic Programming}

The dual formulations of optimal control illuminate different aspects of
transformer behaviour.

\begin{proposition}[Gradient as costate]
During backpropagation through time (BPTT), the gradient of the loss with
respect to the hidden state at time $t$,
\[
  \bm{\lambda}_t = \frac{\partial \cL}{\partial \bm{h}_t},
\]
satisfies the adjoint equation
\begin{equation}\label{eq:adjoint}
  \bm{\lambda}_t
  = \frac{\partial r_t}{\partial \bm{h}_t}
    + \left(\frac{\partial F_\theta}{\partial \bm{h}_t}\right)^\top \bm{\lambda}_{t+1},
\end{equation}
which is the discrete-time Pontryagin costate equation. The gradient
$\bm{\lambda}_t$ decays backward with rate governed by the spectral radius of
$\partial F_\theta / \partial \bm{h}_t$---the same quantity controlling forward
influence decay (\Cref{sec:decay}).
\end{proposition}

\begin{remark}[Symmetry of forward and backward decay]
The contraction rate $\lambda$ in \Cref{thm:decay} governs both:
\begin{enumerate}[label=(\alph*),nosep]
  \item \emph{Forward}: decay of influence of past tokens on future predictions.
  \item \emph{Backward}: decay of gradient signal (vanishing gradients in BPTT).
\end{enumerate}
This is not a coincidence: it is the primal-dual symmetry of optimal control,
where the forward dynamics and backward adjoint share the same Jacobian
structure.
\end{remark}


%=============================================================================
\section{The Bayesian Interpretation}\label{sec:bayes}
%=============================================================================

\subsection{Pretrained Weights as a Prior Measure}

\begin{definition}[Sequence prior]
The pretrained model $p_\theta$ defines a probability measure $\mu_\theta$ over
$\vocab^* = \bigcup_{t=1}^{\infty} \vocab^t$ (the space of all finite token
sequences, where termination occurs at the first
$\langle\mathrm{eos}\rangle$ emission) via the autoregressive factorisation:
\begin{equation}\label{eq:prior}
  \mu_\theta(\bm{x}_{1:\tau})
  = \prod_{t=1}^{\tau} p_\theta(x_t \mid \bm{x}_{1:t-1}),
\end{equation}
where $\tau = \inf\{t : x_t = \langle\mathrm{eos}\rangle\}$ is the stopping
time. We call $\mu_\theta$ the \emph{pretrained prior}. It is a measure over
variable-length sequences with endogenous termination.
\end{definition}

\begin{proposition}[Embedding space as prior geometry]\label{prop:prior_geom}
The token embedding matrix $W_{\mathrm{emb}} \in \R^{d \times V}$ induces a
geometry on $\vocab$ via the inner product
\[
  \langle v, w \rangle_{\mathrm{emb}} = \bm{e}_v^\top \bm{e}_w,
\]
where $\bm{e}_v = W_{\mathrm{emb}}[\cdot, v]$. This geometry encodes the
prior's belief about semantic relatedness: tokens that are substitutable
in many contexts (and hence have similar conditional distributions under
$\mu_\theta$) are embedded nearby. Formally, if
$\KL(p_\theta(\cdot | x_t = v) \| p_\theta(\cdot | x_t = w))$ is small for
most contexts, then $\|\bm{e}_v - \bm{e}_w\|$ tends to be small.
\end{proposition}

This establishes that the embedding space is not merely a lookup table but a
\emph{learned prior geometry} encoding distributional similarity.

\subsection{RLHF as Prior Reweighting}

\begin{theorem}[RLHF posterior]\label{thm:rlhf_bayes}
The optimal RLHF policy~\eqref{eq:gibbs} has a Bayesian interpretation:
\begin{align}
  \underbrace{\pi^*(\bm{x})}_{\text{posterior}}
  &\propto
  \underbrace{\mu_\theta(\bm{x})}_{\text{prior}} \cdot
  \underbrace{\exp\!\bigl(\beta^{-1} R_\phi(\bm{x})\bigr)}_{\text{likelihood (value)}}.
  \label{eq:bayes_rlhf}
\end{align}
The reward model $R_\phi$ plays the role of a log-likelihood, and $\beta$
controls the strength of the evidence relative to the prior. As
$\beta \to \infty$, $\pi^* \to \mu_\theta$ (prior dominates); as
$\beta \to 0$, $\pi^*$ concentrates on $\argmax R_\phi$ (likelihood
dominates).
\end{theorem}

\begin{remark}[Three layers of prior structure]
The full Bayesian hierarchy in a deployed conversational model comprises:
\begin{enumerate}[nosep]
  \item \textbf{Pretrained prior} $\mu_\theta$: broad beliefs about linguistic
        plausibility, learned from the training corpus. Encodes \emph{what can
        follow what}.
  \item \textbf{RLHF-shaped prior} $\pi^*_\theta$: the prior reweighted by
        human values. Encodes \emph{what should follow what} in general.
  \item \textbf{In-context posterior} $\pi^*_\theta(\cdot \mid \bm{c})$: the
        RLHF prior further updated by the conversation context $\bm{c}$.
        Encodes \emph{what should follow what for this user, now}.
\end{enumerate}
\end{remark}

\subsection{In-Context Learning as Amortised Bayesian Inference}

\begin{definition}[In-context posterior]
Given a conversation context $\bm{c} = (c_1, \ldots, c_m)$ comprising the
user's messages and the model's prior responses, the \emph{in-context posterior}
over continuations $\bm{x}_{m+1:T}$ is
\begin{equation}\label{eq:icl_posterior}
  p_\theta(\bm{x}_{m+1:T} \mid \bm{c})
  = \frac{\mu_\theta(\bm{c},\, \bm{x}_{m+1:T})}{\mu_\theta(\bm{c})}.
\end{equation}
\end{definition}

\begin{theorem}[ICL as amortised variational inference]\label{thm:icl}
The transformer's forward pass implements an \emph{amortised} approximation to
Bayesian updating. Specifically, the hidden state $\bm{h}_m = \Phi_\theta(\bm{c})$
serves as a \emph{sufficient statistic} for the posterior over continuations:
\[
  p_\theta(\bm{x}_{m+1:T} \mid \bm{c})
  = p_\theta(\bm{x}_{m+1:T} \mid \bm{h}_m).
\]
The attention mechanism computes the posterior by (approximately) evaluating the
evidence integral
\[
  p_\theta(\bm{c}) = \sum_{\bm{z}} p_\theta(\bm{c} \mid \bm{z})\, p_\theta(\bm{z}),
\]
where $\bm{z}$ ranges over latent ``task'' or ``user'' representations, and the
summation is approximated by the feed-forward network's learned amortisation.
\end{theorem}

\begin{remark}[Conversational Bayesian updating]
In a multi-turn conversation, each user message $c_t$ serves as new evidence
that updates the posterior:
\begin{equation}\label{eq:conv_update}
  p_\theta(\bm{z} \mid c_1, \ldots, c_t)
  \propto p_\theta(c_t \mid \bm{z})\, p_\theta(\bm{z} \mid c_1, \ldots, c_{t-1}).
\end{equation}
The model is inferring a latent representation $\bm{z}$ of the user's intent,
expertise level, communication style, and desired outcomes. This posterior
concentrates over the course of the conversation, explaining why model responses
become more tailored as the dialogue progresses.
\end{remark}

\subsection{Posterior Concentration and Prompt Engineering}

\begin{proposition}[Prompt specificity controls posterior width]\label{prop:prompt}
Let $\bm{c}$ be a prompt and let $\cH_\theta(\bm{c})$ denote the entropy of
the predictive distribution:
\[
  \cH_\theta(\bm{c})
  = -\sum_{v \in \vocab} p_\theta(v \mid \bm{c})\, \log p_\theta(v \mid \bm{c}).
\]
More specific prompts yield lower entropy (tighter posteriors):
\[
  \bm{c} \sqsubseteq \bm{c}' \quad \Longrightarrow \quad
  \E[\cH_\theta(\bm{c}')] \leq \E[\cH_\theta(\bm{c})],
\]
where $\sqsubseteq$ denotes that $\bm{c}'$ is an extension (refinement) of
$\bm{c}$, and the expectation is over the randomness of the extension.
This follows directly from the information-theoretic data processing inequality.
\end{proposition}

\begin{corollary}[Prompt engineering as likelihood design]
Prompt engineering is the craft of constructing a context $\bm{c}$ that serves
as an effective likelihood function: one that concentrates the posterior
$p_\theta(\cdot \mid \bm{c})$ onto the desired region of sequence space.
A vague prompt is a weak likelihood (flat, uninformative); a precise prompt
is a strong likelihood (peaked, informative).
\end{corollary}

%=============================================================================
\section{Unification: The Trinity}\label{sec:unify}
%=============================================================================

We now prove that the three perspectives developed in
Sections~\ref{sec:hmm}--\ref{sec:bayes} are not merely analogies but are
\emph{mathematically identical} descriptions of the same object: the
partition-function ratio $\Gamma(\bm{h}_t, v)$ from~\eqref{eq:Z_ratio}.

\subsection{The Pivotal Object}

Recall from \Cref{thm:opt_policy_Z} that the optimal RLHF policy at each
step takes the form
$\pi^*(v \mid \bm{h}_t)
= \pi_{\mathrm{ref}}(v \mid \bm{h}_t) \cdot \Gamma(\bm{h}_t, v)$
for all $v \in \vocab$ (including $\langle\mathrm{eos}\rangle$), where
$\Gamma$ is defined by~\eqref{eq:Z_ratio}. The following theorem
establishes that $\Gamma$ is simultaneously the operative object in all
three frameworks---and that the endogenous stopping decision is naturally
embedded in each.

\subsection{The Central Equivalence}

\begin{theorem}[The Trinity]\label{thm:trinity}
Let $\pi_{\mathrm{ref}}$ be a pretrained autoregressive policy with hidden
states $\bm{h}_t = \Phi_\theta(\bm{x}_{1:t})$, let
$R_\phi: \vocab^* \to \R$ be a trajectory-level reward defined on
variable-length sequences, and let $\pi^*$ be the optimal policy of the
KL-regularised objective~\eqref{eq:rlhf_obj} with endogenous stopping.
Then the ratio $\Gamma(\bm{h}_t, v)$ from~\eqref{eq:Z_ratio} admits
three equivalent characterisations:

\medskip
\noindent\textbf{(I) Optimal Control: Exponentiated Advantage with Stopping.}

For continuation tokens $v \neq \langle\mathrm{eos}\rangle$:
\begin{equation}\label{eq:trinity_control}
  \Gamma(\bm{h}_t, v)
  = \exp\!\bigl(\beta^{-1} A^*(\bm{h}_t, v)\bigr),
\end{equation}
where the \emph{soft advantage function} is
$A^*(\bm{h}_t, v) = V^*(F_\theta(\bm{h}_t, v)) - V^*(\bm{h}_t)$
and $V^* = \beta \log Z$ satisfies the free-boundary soft Bellman
equation~\eqref{eq:soft_bellman}.

For the stopping action:
\begin{equation}\label{eq:trinity_control_eos}
  \Gamma(\bm{h}_t, \langle\mathrm{eos}\rangle)
  = \exp\!\bigl(\beta^{-1} [R_\phi(\bm{x}_{1:t}) - V^*(\bm{h}_t)]\bigr).
\end{equation}
The model stops when the immediate reward $R_\phi(\bm{x}_{1:t})$ exceeds
the value-to-go $V^*(\bm{h}_t)$: the option is exercised when the
current payoff dominates the expected future payoff.

\medskip
\noindent\textbf{(II) Bayesian Inference: Sequential Likelihood Ratio with
Stopping as Commitment.}
\begin{equation}\label{eq:trinity_bayes}
  \Gamma(\bm{h}_t, v)
  = \frac{\pi^*(v \mid \bm{h}_t)}{\pi_{\mathrm{ref}}(v \mid \bm{h}_t)}
  = \frac{\text{posterior predictive}}{\text{prior predictive}},
\end{equation}
for all $v \in \vocab$ (including $\langle\mathrm{eos}\rangle$). For
continuation tokens, $\Gamma$ measures how much choosing $v$ shifts the
marginal likelihood of producing a high-value trajectory. For
$v = \langle\mathrm{eos}\rangle$, $\Gamma$ measures the posterior
probability that the \emph{evidence is sufficient}: stopping corresponds
to the Bayesian decision that no further tokens would improve the expected
reward. This is the \emph{sequential probability ratio test} (SPRT) of
Wald: the model continues gathering ``evidence'' (emitting tokens) until
the posterior concentrates enough to commit.

\medskip
\noindent\textbf{(III) Hidden Markov Model: Emission-Reweighted Transition
with Absorbing State.}
\begin{equation}\label{eq:trinity_hmm}
  \Gamma(\bm{h}_t, v)
  = \frac{\kappa^*\bigl((\bm{h}_t, x_t),\, F_\theta(\bm{h}_t, v) \times \{v\}\bigr)}%
         {\kappa_{\mathrm{ref}}\bigl((\bm{h}_t, x_t),\, F_\theta(\bm{h}_t, v) \times \{v\}\bigr)},
\end{equation}
where $\kappa^*$ is the transition kernel of the hidden Markov process under
$\pi^*$ and $\kappa_{\mathrm{ref}}$ is the kernel under $\pi_{\mathrm{ref}}$.
For $v = \langle\mathrm{eos}\rangle$, the transition leads to the
\emph{absorbing state}: the HMM is killed. The optimal HMM reweights
the absorption probability relative to the reference HMM by
$\Gamma(\bm{h}_t, \langle\mathrm{eos}\rangle)$, increasing it when
stopping is advantageous and decreasing it when continuation is preferred.
\end{theorem}

\begin{proof}
We prove each characterisation and their mutual equivalence.

\medskip
\noindent\textbf{Proof of (I).}
For continuation tokens $v \neq \langle\mathrm{eos}\rangle$,
by definition~\eqref{eq:soft_V}:
\[
  \Gamma(\bm{h}_t, v)
  = \frac{Z(F_\theta(\bm{h}_t, v))}{Z(\bm{h}_t)}
  = \frac{\exp(\beta^{-1} V^*(F_\theta(\bm{h}_t, v)))}%
         {\exp(\beta^{-1} V^*(\bm{h}_t))}
  = \exp\!\bigl(\beta^{-1} A^*(\bm{h}_t, v)\bigr).
\]
For $v = \langle\mathrm{eos}\rangle$:
\[
  \Gamma(\bm{h}_t, \langle\mathrm{eos}\rangle)
  = \frac{\exp(\beta^{-1} R_\phi(\bm{x}_{1:t}))}%
         {\exp(\beta^{-1} V^*(\bm{h}_t))}
  = \exp\!\bigl(\beta^{-1}[R_\phi(\bm{x}_{1:t}) - V^*(\bm{h}_t)]\bigr).
\]
The stopping advantage $R_\phi(\bm{x}_{1:t}) - V^*(\bm{h}_t)$ is positive
precisely on the optimal stopping region $\cS^*$: the model stops when the
immediate reward exceeds the continuation value. The Bellman
decomposition---including the stopping boundary---is forced by the structure
of the problem: trajectory-level reward with sequential token emission and
endogenous termination. \qed

\medskip
\noindent\textbf{Proof of (II).}
From \Cref{thm:opt_policy_Z},
$\pi^*(v \mid \bm{h}_t) = \pi_{\mathrm{ref}}(v \mid \bm{h}_t) \cdot
\Gamma(\bm{h}_t, v)$ for all $v \in \vocab$, so
\[
  \Gamma(\bm{h}_t, v)
  = \frac{\pi^*(v \mid \bm{h}_t)}{\pi_{\mathrm{ref}}(v \mid \bm{h}_t)}.
\]
Now, $\pi^*(\bm{x}_{1:\tau}) \propto \pi_{\mathrm{ref}}(\bm{x}_{1:\tau})
\exp(\beta^{-1} R_\phi(\bm{x}_{1:\tau}))$ is a Bayesian posterior with
prior $\pi_{\mathrm{ref}}$ and log-likelihood $\beta^{-1} R_\phi$
(\Cref{thm:rlhf_bayes}).

For continuation tokens, $\Gamma$ is the per-step likelihood ratio:
how much choosing $v$ updates the belief that the trajectory will be
high-value. For $v = \langle\mathrm{eos}\rangle$, $\Gamma$ is the
\emph{commitment ratio}: it measures the posterior belief that the
accumulated evidence (tokens emitted so far) is sufficient to achieve
high reward without further generation. This is precisely the stopping
rule of the \emph{sequential probability ratio test} (Wald, 1945):
continue sampling (emitting tokens) until the accumulated evidence
crosses a threshold, then commit (emit EOS).

The sequential factorisation extends to variable-length sequences:
\begin{equation}\label{eq:sequential_bayes}
  \frac{\pi^*(\bm{x}_{1:\tau})}{\pi_{\mathrm{ref}}(\bm{x}_{1:\tau})}
  = \prod_{t=1}^{\tau} \Gamma(\bm{h}_{t-1}, x_t),
\end{equation}
where the final factor $\Gamma(\bm{h}_{\tau-1}, \langle\mathrm{eos}\rangle)$
encodes the stopping decision. The trajectory-level Bayes factor decomposes
as a product of per-step likelihood ratios, with the number of factors
itself being a random variable determined by the stopping rule. \qed

\medskip
\noindent\textbf{Proof of (III).}
The transition kernel of the T-HMM (\Cref{def:thmm}) under policy $\pi$
is (cf.~\eqref{eq:kernel}):
\[
  \kappa_\pi\!\bigl((\bm{h}_t, x_t),\, B \times \{v\}\bigr)
  = \pi(v \mid \bm{h}_t) \cdot \ind\!\bigl[F_\theta(\bm{h}_t, v) \in B\bigr].
\]
For $v \neq \langle\mathrm{eos}\rangle$, the state transition $F_\theta$ is
deterministic given $(\bm{h}_t, v)$, and the Radon-Nikod\'ym derivative
reduces to the emission ratio:
\[
  \frac{d\kappa^*}{d\kappa_{\mathrm{ref}}}
  \bigg|_{(\bm{h}_t, x_t),\, v}
  = \frac{\pi^*(v \mid \bm{h}_t)}{\pi_{\mathrm{ref}}(v \mid \bm{h}_t)}
  = \Gamma(\bm{h}_t, v).
\]
For $v = \langle\mathrm{eos}\rangle$, the transition leads to an
\emph{absorbing state} $\bm{h}_\dagger$. The Radon-Nikod\'ym derivative of
the absorption probability is
\[
  \frac{\pi^*(\langle\mathrm{eos}\rangle \mid \bm{h}_t)}%
       {\pi_{\mathrm{ref}}(\langle\mathrm{eos}\rangle \mid \bm{h}_t)}
  = \Gamma(\bm{h}_t, \langle\mathrm{eos}\rangle).
\]
The optimal HMM is obtained from the reference HMM by tilting
\emph{both} the continuation emissions and the absorption probability
by $\Gamma$. \qed

\medskip
\noindent\textbf{Mutual consistency.}
The three characterisations are three readings of the single identity,
valid for all $v \in \vocab$ including $\langle\mathrm{eos}\rangle$:
\begin{equation}\label{eq:master}
  \boxed{%
    \underbrace{\exp\!\bigl(\beta^{-1} A^*(\bm{h}_t, v)\bigr)}_{\text{(I) advantage / stopping value}}
    \;=\;
    \underbrace{\frac{\pi^*(v|\bm{h}_t)}{\pi_{\mathrm{ref}}(v|\bm{h}_t)}}_{\text{(II) likelihood ratio / SPRT}}
    \;=\;
    \underbrace{\frac{d\kappa^*}{d\kappa_{\mathrm{ref}}}
    \bigg|_{(\bm{h}_t, x_t),\, v}}_{\text{(III) kernel tilt / absorption}}
    \;=\;
    \Gamma(\bm{h}_t, v).
  }
\end{equation}
All three are \emph{defined} as the same ratio from~\eqref{eq:Z_ratio}
and therefore agree identically, not approximately. When $v =
\langle\mathrm{eos}\rangle$, the identity connects the optimal stopping
decision across all three frameworks: the exercise boundary (control),
the commitment threshold (Bayesian), and the absorption rate (HMM) are
the same surface in hidden-state space. \qed
\end{proof}

\subsection{The Necessity of Value Encoding}

\begin{theorem}[Structural necessity of Bellman decomposition]\label{thm:necessity}
Let $\pi_\theta$ be any autoregressive policy that approximates the optimal
RLHF policy $\pi^*$ in the sense that
$\KL(\pi^* \| \pi_\theta) \leq \varepsilon$.
Then there exists a function $\hat{V}: \R^d \to \R$ such that:
\begin{enumerate}[label=(\roman*)]
  \item $\hat{V}$ is computable from the hidden state: $\hat{V}(\bm{h}_t) =
        \psi(\bm{h}_t)$ for some function $\psi$ extractable from the model's
        parameters.
  \item $\hat{V}$ satisfies the soft Bellman equation to within
        $O(\sqrt{\varepsilon})$:
        \[
          \left|\hat{V}(\bm{h}_t) - \beta \log
          \E_{v \sim \pi_{\mathrm{ref}}}\!\left[
            \exp\!\bigl(\beta^{-1} \hat{V}(F_\theta(\bm{h}_t, v))\bigr)
          \right]\right|
          \leq O(\sqrt{\varepsilon}).
        \]
\end{enumerate}
\end{theorem}

\begin{proof}
Define $\hat{V}(\bm{h}_t) = \beta \log \sum_{v \in \vocab}
\pi_\theta(v|\bm{h}_t) / \pi_{\mathrm{ref}}(v|\bm{h}_t) \cdot
\pi_{\mathrm{ref}}(v|\bm{h}_t) \cdot
\exp(\beta^{-1} \hat{V}(F_\theta(\bm{h}_t, v)))$---i.e., the soft value
function \emph{induced by} $\pi_\theta$ as if it were optimal. Since
$\pi_\theta \approx \pi^*$ (in KL), we have
$\pi_\theta(v|\bm{h}_t)/\pi_{\mathrm{ref}}(v|\bm{h}_t) \approx
\Gamma(\bm{h}_t, v)$ pointwise in a total-variation sense (by Pinsker's
inequality, $\|\pi_\theta(\cdot|\bm{h}_t) - \pi^*(\cdot|\bm{h}_t)\|_{\TV}
\leq \sqrt{\varepsilon/2}$). Substituting into the soft Bellman equation
and bounding the residual via the Lipschitz continuity of $\log$ on $[p_{\min}, 1]$
yields the $O(\sqrt{\varepsilon})$ bound. \qed
\end{proof}

\begin{remark}[Interpretation]
\Cref{thm:necessity} says: you cannot be a good RLHF policy without
approximately solving the Bellman equation. The Bellman structure is not
an optional interpretation---it is a \emph{necessary consequence} of
optimising a trajectory-level reward through sequential token emission.
The model has no choice but to encode an approximate value function.
\end{remark}

\subsection{Decomposition of the Global Bayes Factor}

\begin{corollary}[Telescoping identity with endogenous stopping]\label{cor:telescope}
The trajectory-level Bayes factor (posterior-to-prior ratio) decomposes as a
telescoping product of per-step $\Gamma$-ratios, including the stopping
decision:
\begin{equation}\label{eq:telescope}
  \frac{\pi^*(\bm{x}_{1:\tau})}{\pi_{\mathrm{ref}}(\bm{x}_{1:\tau})}
  = \frac{\exp(\beta^{-1} R_\phi(\bm{x}_{1:\tau}))}{Z_0}
  = \prod_{t=1}^{\tau} \Gamma(\bm{h}_{t-1}, x_t)
  = \exp\!\left(\beta^{-1} \sum_{t=1}^{\tau} A^*(\bm{h}_{t-1}, x_t)\right),
\end{equation}
where $x_\tau = \langle\mathrm{eos}\rangle$ and the final factor
$\Gamma(\bm{h}_{\tau-1}, \langle\mathrm{eos}\rangle)$ encodes the
stopping decision.
This identity says three things simultaneously:
\begin{enumerate}[label=(\alph*),nosep]
  \item \textbf{Control:} The trajectory reward decomposes into a sum of
        per-step advantages \emph{including the stopping advantage}---this
        is the reward decomposition lemma of the free-boundary Bellman
        framework. The number of terms is itself optimised.
  \item \textbf{Bayesian:} The full Bayes factor decomposes into a product of
        per-step likelihood ratios, with the final factor being the
        \emph{commitment decision}---the chain rule of Bayesian
        updating with an endogenous sample size (SPRT).
  \item \textbf{HMM:} The Radon-Nikod\'ym derivative of the trajectory
        measure decomposes into per-step kernel ratios terminating at
        absorption---the Girsanov-type change of measure for Markov
        chains with killing.
\end{enumerate}
These three decompositions are the same identity read in three languages.
The endogenous stopping time $\tau$ is the same in all three: the
exercise boundary, the commitment threshold, and the absorption time
are one object.
\end{corollary}

\begin{proof}
For the telescoping product, noting that for
$t < \tau$, $x_t \neq \langle\mathrm{eos}\rangle$ so
$\Gamma(\bm{h}_{t-1}, x_t) = Z(\bm{h}_t)/Z(\bm{h}_{t-1})$, and for
$t = \tau$, $x_\tau = \langle\mathrm{eos}\rangle$ so
$\Gamma(\bm{h}_{\tau-1}, \langle\mathrm{eos}\rangle)
= \exp(\beta^{-1} R_\phi(\bm{x}_{1:\tau-1}))/Z(\bm{h}_{\tau-1})$:
\begin{align}
  \prod_{t=1}^{\tau} \Gamma(\bm{h}_{t-1}, x_t)
  &= \left(\prod_{t=1}^{\tau-1} \frac{Z(\bm{h}_t)}{Z(\bm{h}_{t-1})}\right)
     \cdot \frac{\exp(\beta^{-1} R_\phi(\bm{x}_{1:\tau-1}))}{Z(\bm{h}_{\tau-1})} \notag \\
  &= \frac{Z(\bm{h}_{\tau-1})}{Z(\bm{h}_0)}
     \cdot \frac{\exp(\beta^{-1} R_\phi(\bm{x}_{1:\tau-1}))}{Z(\bm{h}_{\tau-1})}
  = \frac{\exp(\beta^{-1} R_\phi(\bm{x}_{1:\tau}))}{Z_0}.
\end{align}
The telescoping cancellation is exact, and the stopping factor merges
seamlessly. \qed
\end{proof}

\subsection{Summary: Why the Trinity Holds}

The Trinity is not a coincidence, a metaphor, or an approximation. It holds
because:

\begin{enumerate}[nosep]
  \item \textbf{Sequential emission + trajectory reward + endogenous stopping
        $\Rightarrow$ free-boundary Bellman.}
        Any agent that must act step-by-step, is evaluated on the
        cumulative outcome, and decides when to stop is solving a
        free-boundary dynamic programming problem. The soft Bellman
        equation with stopping (\Cref{thm:soft_bellman}) is the
        \emph{unique} decomposition of the trajectory-level RLHF objective
        into per-step decisions plus a stopping rule.

  \item \textbf{Gibbs form + chain rule + endogenous sample size
        $\Rightarrow$ sequential testing.}
        The optimal RLHF policy has the Gibbs form
        (\Cref{thm:rlhf_optimal}), which IS Bayes' rule. Its autoregressive
        factorisation IS sequential Bayesian updating. The endogenous
        stopping time makes this a sequential probability ratio test:
        the model continues until the evidence is sufficient.
        The partition function $Z(\bm{h}_t)$ IS the marginal likelihood.

  \item \textbf{Markov state + deterministic transitions + absorbing state
        $\Rightarrow$ killed HMM.}
        The hidden state is Markov by construction (\Cref{prop:markov}).
        The EOS token induces an absorbing state, creating a Markov chain
        with killing. The change from reference to optimal policy is a
        Radon-Nikod\'ym tilt of both the emission kernel and the
        absorption rate.

  \item \textbf{All three meet at $\Gamma$, including at the boundary.}
        The ratio $\Gamma$ is the unique object that simultaneously
        encodes the per-step advantage, the per-step likelihood ratio,
        and the per-step kernel derivative. For $v \neq
        \langle\mathrm{eos}\rangle$, it governs continuation. For
        $v = \langle\mathrm{eos}\rangle$, it governs stopping. The
        optimal stopping region---the exercise boundary, the commitment
        threshold, the absorption surface---is the same set $\cS^*$ in
        all three frameworks.
\end{enumerate}

\subsection{Convergence of the Three Decay Mechanisms}

A further consistency check: all three frameworks predict the same memory
decay, governed by the same spectral parameter.

\begin{proposition}[Unified decay]\label{prop:unified_decay}
Under the contraction condition of \Cref{thm:decay}:
\begin{enumerate}[label=(\alph*),nosep]
  \item \textbf{HMM}: Influence of token $x_s$ on prediction at $t$ decays as
        $O(\lambda^{2(t-s)})$ (forward mixing).
  \item \textbf{Control}: The advantage $A^*(\bm{h}_t, v)$'s sensitivity to a
        perturbation at step $s$ decays as $O(\lambda^{t-s})$, since
        $A^*$ depends on $V^*(\bm{h}_t)$ which inherits the contraction of
        $\bm{h}_t$ (cf.~\eqref{eq:adjoint}).
  \item \textbf{Bayesian}: The per-step likelihood ratio $\Gamma(\bm{h}_t, v)$'s
        sensitivity to evidence at step $s$ decays as $O(\lambda^{2(t-s)})$
        (posterior decorrelation).
\end{enumerate}
All three rates are governed by the same quantity $\lambda$: the contraction
rate of the hidden-state dynamics $F_\theta$.
\end{proposition}


%=============================================================================
\section{Implications and Interpretive Consequences}\label{sec:implications}
%=============================================================================

\subsection{Coherence Degradation in Long Generations}

\begin{proposition}[Ergodic drift]
For generation lengths $T \gg \tau_{\mathrm{mix}}$, the autoregressive process
converges toward the stationary distribution of the Markov chain, losing
coherence with the initial prompt. This is the long-run memorylessness described
in \Cref{cor:memoryless}: the ``plot'' set up in the opening tokens is
gradually forgotten as the process mixes.
\end{proposition}

This explains the empirical observation that very long model generations tend to
become generic, repetitive, or drift from the original topic---even when the
context window has not been exhausted.

\subsection{Chain-of-Thought as Control Horizon Extension}

In the control framework, CoT prompting expands the effective planning horizon.

\begin{proposition}
Let $T_0$ be the base horizon and $T_{\mathrm{CoT}} = T_0 + K$ the
horizon with $K$ additional reasoning tokens. The value achievable satisfies
\[
  V^*_{T_{\mathrm{CoT}}}(\bm{h}_0) \geq V^*_{T_0}(\bm{h}_0)
\]
with strict inequality whenever the optimal trajectory under $T_0$ is
suboptimal---i.e., whenever intermediate reasoning steps allow the controller
to reach higher-value terminal states.
\end{proposition}

\subsection{Few-Shot Learning as Strong Likelihood}

In the Bayesian framework, few-shot examples act as concentrated evidence.

\begin{proposition}
Let $\bm{c}_k$ be a prompt containing $k$ input-output demonstrations. The
posterior entropy satisfies
\[
  \cH_\theta(\bm{c}_k) \leq \cH_\theta(\bm{c}_{k-1}) \leq \cdots \leq
  \cH_\theta(\bm{c}_0),
\]
with equality iff the additional examples are redundant given the
existing context. Each example narrows the posterior, concentrating the
model's predictive distribution on the task-relevant region of sequence space.
\end{proposition}

\subsection{Temperature as Posterior Sharpening}

\begin{proposition}
Sampling with temperature $\tau > 0$ modifies the emission distribution to
\[
  p_\tau(v \mid \bm{h}_t)
  = \frac{p_\theta(v \mid \bm{h}_t)^{1/\tau}}%
         {\sum_{w} p_\theta(w \mid \bm{h}_t)^{1/\tau}}.
\]
In the Bayesian framework, $\tau < 1$ corresponds to \emph{tempering} the
posterior (increasing confidence beyond what the evidence warrants), while
$\tau > 1$ corresponds to \emph{flattening} (increased epistemic humility).
In the control framework, $\tau$ modulates the \emph{exploration-exploitation
trade-off} of the stochastic policy.
\end{proposition}


%=============================================================================
\section{Formal Connections to Classical Theory}\label{sec:connections}
%=============================================================================

\subsection{Relationship to Predictive State Representations}

The T-HMM of \Cref{def:thmm} is closely related to Predictive State
Representations (PSRs). In a PSR, the state is defined as a vector of
predictions about future observables, and the system is ``observable'' in the
sense that the state can be maintained from observations alone.

\begin{proposition}
The transformer hidden state $\bm{h}_t$ constitutes an \emph{approximate}
predictive state: it encodes (approximately) all information about the
future predictive distribution that is extractable from the past observations
$\bm{x}_{1:t}$. It differs from a classical PSR in that:
\begin{enumerate}[label=(\alph*),nosep]
  \item The state is computed by a learned, parametric function
        $\Phi_\theta$ rather than maintained via an update rule derived
        from the system's dynamics.
  \item The approximation quality depends on model capacity and training
        adequacy.
\end{enumerate}
\end{proposition}

\subsection{Relationship to Information Bottleneck}

\begin{proposition}
The hidden state $\bm{h}_t$ solves an approximate \emph{information bottleneck}
problem:
\[
  \min_{\bm{h}_t}\; I(\bm{x}_{1:t};\, \bm{h}_t) - \beta_{\mathrm{IB}}\, I(\bm{h}_t;\, x_{t+1}),
\]
where $I(\cdot;\cdot)$ denotes mutual information. The state compresses the
past ($\bm{x}_{1:t}$) while retaining information predictive of the future
($x_{t+1}$). The contraction rate $\lambda$ controls how aggressively the
bottleneck compresses distant information.
\end{proposition}

\subsection{Relationship to Free Energy Principle}

The RLHF objective~\eqref{eq:rlhf_obj} has a direct connection to variational
free energy minimisation:
\begin{align}
  -J_{\mathrm{RLHF}}(\theta)
  &= -\E_{\pi_\theta}[R_\phi] + \beta\, \KL(\pi_\theta \| \pi_{\mathrm{ref}}) \\
  &= \beta\!\left(
       -\beta^{-1}\E_{\pi_\theta}[R_\phi]
       + \KL(\pi_\theta \| \pi_{\mathrm{ref}})
     \right) \\
  &= \beta\, \underbrace{\bigl(
       \E_{\pi_\theta}[-\log \tilde{p}(\bm{x})]
       + \KL(\pi_\theta \| \pi_{\mathrm{ref}})
     \bigr)}_{\text{variational free energy}},
\end{align}
where $\tilde{p}(\bm{x}) \propto \exp(\beta^{-1} R_\phi(\bm{x}))$. Minimising
$-J_{\mathrm{RLHF}}$ is thus equivalent to minimising the variational free
energy, connecting transformer training to the free energy principle from
theoretical neuroscience.


%=============================================================================
\section{The Spectral Structure of Memory}\label{sec:spectral}
%=============================================================================

To make the decay analysis of \Cref{sec:decay} more precise, we study the
spectral properties of the linearised hidden-state dynamics.

\subsection{Transfer Operator}

\begin{definition}[Linearised transfer operator]
At a reference trajectory $\bar{\bm{h}}_{1:T}$, define the \emph{transfer
operator} at step $t$ as
\[
  \cT_t = \frac{\partial F_\theta(\bar{\bm{h}}_t, \bar{x}_{t+1})}%
               {\partial \bm{h}_t}
  \in \R^{d \times d}.
\]
The cumulative transfer from step $s$ to step $t$ is
\[
  \cT_{s:t} = \prod_{\tau=s}^{t-1} \cT_\tau.
\]
\end{definition}

\begin{theorem}[Spectral characterisation of memory]\label{thm:spectral}
Let $\sigma_1(\cT_{s:t}) \geq \cdots \geq \sigma_d(\cT_{s:t})$ be the singular
values of $\cT_{s:t}$. Then:
\begin{enumerate}[label=(\roman*)]
  \item The \emph{effective rank} of the influence propagation from $s$ to $t$ is
  \[
    r_{\mathrm{eff}}(s,t)
    = \frac{\bigl(\sum_i \sigma_i\bigr)^2}{\sum_i \sigma_i^2}
    \leq d,
  \]
  and this decreases with $|t-s|$ as modes decay at different rates.

  \item The \emph{memory spectrum} $\{\lambda_k\}_{k=1}^d$, defined as the
  geometric means $\lambda_k = (\prod_\tau \sigma_k(\cT_\tau))^{1/(t-s)}$,
  determines which directions in hidden space carry long-range information.

  \item Directions with $\lambda_k \approx 1$ correspond to
  \emph{persistent memory modes} (e.g., topic, style, factual anchors),
  while directions with $\lambda_k \ll 1$ carry only local information
  (e.g., recent syntactic context).
\end{enumerate}
\end{theorem}

This provides a mechanistic explanation for why transformers can simultaneously
exhibit short-range syntactic sensitivity and long-range semantic coherence: the
memory is \emph{spectrally structured}, with different components decaying at
different rates.


%=============================================================================
\section{Discussion and Open Questions}\label{sec:discussion}
%=============================================================================

\subsection{Time Inhomogeneity}

The T-HMM is fundamentally \emph{time-inhomogeneous}: the effective transition
kernel $\kappa_\theta$ varies with position (due to positional encodings) and
with context (due to attention). Classical ergodic theory for time-homogeneous
chains must be extended to handle this case. The relevant framework is that of
\emph{non-stationary Markov chains} with \emph{uniform ergodicity} conditions.

\subsection{Non-Monotonic Attention}

The contraction assumption in \Cref{thm:decay} is violated when attention heads
``reach back'' to specific positions with high weight. This creates
\emph{punctuated memory}: the influence function $I(s,t)$ is not monotonically
decreasing but exhibits spikes at positions of high attention. A more refined
model would replace the uniform contraction with a \emph{random} contraction
rate, leading to a multiplicative random walk analysis.

\subsection{Scaling Laws and the Prior}

As model scale increases, the pretrained prior $\mu_\theta$ becomes a
better approximation to the ``true'' distribution over natural language.
In the Bayesian framework, this means that larger models require \emph{less}
evidence (shorter prompts) to concentrate on the correct posterior---which
is consistent with the empirical finding that larger models exhibit stronger
few-shot learning.

\subsection{Open Problems}

\begin{enumerate}[nosep]
  \item \textbf{Exact characterisation of the spectral gap.} Can the mixing
        time of the T-HMM be bounded in terms of architectural parameters
        (depth, width, number of heads)?

  \item \textbf{Finite-sample Bayesian consistency.} Under what conditions does
        the in-context posterior converge to the true task posterior as the number
        of in-context examples grows?

  \item \textbf{Control complexity.} What is the computational complexity of the
        optimal control problem solved (approximately) by the transformer, and
        how does it relate to the model's inference cost?

  \item \textbf{Value function geometry.} Can the implicit value function
        $V_\theta(\bm{h}_t)$ be extracted and analysed? Does it exhibit the
        smoothness properties assumed in approximate dynamic programming?

  \item \textbf{Multi-agent control.} In conversational settings, the user and
        the model are both acting as controllers with coupled dynamics. The
        appropriate framework may be \emph{Stackelberg games} or
        \emph{cooperative control} rather than single-agent optimisation.
\end{enumerate}

%=============================================================================
\section{Conclusion}\label{sec:conclusion}
%=============================================================================

We have developed a unified theoretical framework that characterises
autoregressive transformer language models through three equivalent lenses:

\begin{enumerate}[nosep]
  \item As \textbf{hidden Markov processes} with continuous-state dynamics and
        exponentially decaying memory, formalising the notion that these models
        are ``memoryless in the long run'' while retaining local, vanishingly
        influential memory at each step.

  \item As \textbf{stochastic optimal controllers} that solve a Bellman equation
        over the hidden-state manifold, selecting tokens as control actions that
        steer the sequence toward high-value terminal states.

  \item As \textbf{Bayesian inference engines} whose pretrained weights define a
        prior over language, RLHF reshapes this prior with a value-based
        likelihood, and in-context conditioning performs amortised posterior
        updating.
\end{enumerate}

The three frameworks are not merely analogies but are provably equivalent
formalisations, with the hidden state $\bm{h}_t$ playing the unifying role
of Markov state, control state, and sufficient statistic simultaneously.
The spectral structure of the hidden-state dynamics governs the decay of
influence in all three frameworks, providing a single quantitative
characterisation of the model's effective memory.

This unified perspective yields interpretive power: chain-of-thought prompting
is control horizon extension; prompt engineering is likelihood design; RLHF is
prior reweighting; temperature is posterior tempering; coherence degradation in
long generations is ergodic mixing. These are not metaphors---they are
mathematical consequences of the framework.

\bigskip

\appendix

%=============================================================================
\section{Detailed Proof of \Cref{thm:decay}}\label{app:decay_proof}
%=============================================================================

\begin{proof}
Let $\bm{x}_{1:t}$ be a token sequence and let $\bm{x}_{1:t}^{\setminus s}$
denote the sequence with $x_s$ replaced by an independent draw $\tilde{x}_s
\sim p_\theta(x_s)$. Define
\[
  \bm{h}_\tau = \Phi_\theta(\bm{x}_{1:\tau}), \qquad
  \tilde{\bm{h}}_\tau = \Phi_\theta(\bm{x}_{1:\tau}^{\setminus s}).
\]
For $\tau < s$, $\bm{h}_\tau = \tilde{\bm{h}}_\tau$. At $\tau = s$, the
perturbation introduces a discrepancy
$\|\bm{h}_s - \tilde{\bm{h}}_s\| \leq D$
where $D$ is the diameter of the embedding space.

For $\tau > s$, since $\bm{h}_{\tau+1} = F_\theta(\bm{h}_\tau, x_{\tau+1})$
and $F_\theta$ is $\lambda$-Lipschitz in its first argument:
\[
  \|\bm{h}_\tau - \tilde{\bm{h}}_\tau\| \leq \lambda^{\tau - s}\, D.
\]

The emission distributions at step $t$ are
\[
  p_\theta(\cdot \mid \bm{h}_t), \qquad
  p_\theta(\cdot \mid \tilde{\bm{h}}_t).
\]
Since $g_\theta = \softmax \circ (W_{\mathrm{out}} \cdot)$ is Lipschitz with
constant $L_g$ (the softmax Jacobian is bounded by $\nicefrac{1}{4}$ per
coordinate, and $W_{\mathrm{out}}$ has bounded operator norm),
\[
  \|g_\theta(\bm{h}_t) - g_\theta(\tilde{\bm{h}}_t)\|_1
  \leq L_g\, \|\bm{h}_t - \tilde{\bm{h}}_t\|
  \leq L_g\, D\, \lambda^{t-s}.
\]

By Pinsker's inequality and the local quadratic bound on KL divergence:
\[
  \KL\!\bigl(g_\theta(\bm{h}_t) \| g_\theta(\tilde{\bm{h}}_t)\bigr)
  \leq \frac{\|g_\theta(\bm{h}_t) - g_\theta(\tilde{\bm{h}}_t)\|_1^2}%
            {2\, \min_v p_\theta(v|\tilde{\bm{h}}_t)}
  \leq C'\, \lambda^{2(t-s)},
\]
where $C' = L_g^2 D^2 / (2 p_{\min})$ and $p_{\min}$ is a lower bound on
the emission probabilities (which exists for softmax with bounded logits).

Taking expectations over the randomness of $\tilde{x}_s$:
\[
  I(s,t) = \E\!\left[\KL\!\bigl(g_\theta(\bm{h}_t) \| g_\theta(\tilde{\bm{h}}_t)\bigr)\right]
  \leq C'\, \lambda^{2(t-s)} \eqqcolon C\, \lambda^{2(t-s)}.
  \qedhere
\]
\end{proof}


%=============================================================================
\section{Proof of \Cref{thm:rlhf_optimal}}\label{app:rlhf_proof}
%=============================================================================

\begin{proof}
The RLHF objective is
\[
  J(\pi) = \E_\pi[R_\phi(\bm{x})] - \beta\, \KL(\pi \| \pi_{\mathrm{ref}})
  = \sum_{\bm{x}} \pi(\bm{x})\!\left[
    R_\phi(\bm{x}) - \beta \log \frac{\pi(\bm{x})}{\pi_{\mathrm{ref}}(\bm{x})}
  \right].
\]
Setting the functional derivative to zero:
\[
  \frac{\delta J}{\delta \pi(\bm{x})}
  = R_\phi(\bm{x}) - \beta \log \frac{\pi(\bm{x})}{\pi_{\mathrm{ref}}(\bm{x})} - \beta
  = 0.
\]
Solving:
\[
  \log \pi^*(\bm{x})
  = \log \pi_{\mathrm{ref}}(\bm{x}) + \beta^{-1} R_\phi(\bm{x}) - 1.
\]
Exponentiating:
\[
  \pi^*(\bm{x}) \propto \pi_{\mathrm{ref}}(\bm{x})\, \exp\!\bigl(\beta^{-1} R_\phi(\bm{x})\bigr).
\]
The proportionality constant is the normalising partition function
$Z = \sum_{\bm{x}} \pi_{\mathrm{ref}}(\bm{x})\, \exp(\beta^{-1} R_\phi(\bm{x}))$.
Strict concavity of $J$ in $\pi$ (since $-\KL$ is strictly concave)
guarantees uniqueness. \qed
\end{proof}


%=============================================================================
\section{Notation Reference}\label{app:notation}
%=============================================================================

\begin{center}
\renewcommand{\arraystretch}{1.3}
\begin{tabular}{@{}lp{10cm}@{}}
  \toprule
  \textbf{Symbol} & \textbf{Meaning} \\
  \midrule
  $\vocab$ & Token vocabulary, $|\vocab| = V$ \\
  $d$ & Hidden state dimension \\
  $T$ & Sequence length / control horizon \\
  $\bm{x}_{1:t}$ & Token sequence $(x_1, \ldots, x_t)$ \\
  $\bm{h}_t$ & Hidden state at step $t$, $\bm{h}_t \in \R^d$ \\
  $\Phi_\theta$ & Transformer encoding map \\
  $F_\theta$ & Incremental hidden-state update \\
  $g_\theta$ & Emission distribution (softmax output) \\
  $\kappa_\theta$ & Markov transition kernel \\
  $\mu_\theta$ & Pretrained prior over sequences \\
  $\pi_\theta$ & Policy (conditional token distribution) \\
  $R_\phi$ & Learned reward model \\
  $V_\theta$ & Value function (value-to-go) \\
  $\lambda$ & Contraction rate of hidden dynamics \\
  $\beta$ & KL penalty / inverse temperature \\
  $\tau$ & Sampling temperature \\
  $\tau_{\mathrm{mix}}$ & Mixing time of the hidden Markov process \\
  \bottomrule
\end{tabular}
\end{center}

\end{document}
